
\documentclass[a4paper,12pt]{article}

% --- Codificação e idioma ---
\usepackage[utf8]{inputenc}
\usepackage[T1]{fontenc}
\usepackage[portuguese]{babel}

% --- Formatação ---
\usepackage[a4paper,margin=2.5cm]{geometry}
\usepackage{lmodern}
\usepackage{microtype}
\usepackage{setspace}
\usepackage{parskip}

% --- Links ---
\usepackage[hidelinks]{hyperref}

% --- Dados do documento ---
\title{ AeroEye}
\author{
    Daniel Santana, Gabriel Bernardes, Reiben Alves \\
    \small Desenvolvimento de Dispositivos Móveis \\
    \small Universidade de Leiria
}
\date{}

\begin{document}

\maketitle

\section{Contextualização do projeto}

O AeroEye consiste num sistema de monitorização de maquinaria industrial baseado num drone, que recorre a técnicas de Machine Learning para identificar potenciais anomalias e a tecnologias de Realidade Virtual (VR) e Realidade Aumentada (AR) para facilitar o acesso remoto de um técnico especializado à anomalia e fornecer apoio visual à manutenção. 

O projeto pretende reduzir os custos operacionais e as paragens de produção, evitando deslocações desnecessárias de técnicos externos quando o problema não é de dificil resolução e permitindo que os operadores resolvam determinados problemas mais rapidamente com as devidas instruções.

\section{Desenvolvimento no âmbito de DDM}
No âmbito da unidade curricular de Desenvolvimento de Dispositivos Móveis, será desenvolvida uma aplicação nativa para IoS, integrada no sistema AeroEye. A aplicação permitirá consultar relatórios e vídeos, bem como visualizar, em tempo real, as imagens das câmaras e os dados recolhidos pelos sensores do drone. Adicionalmente, incluirá a funcionalidade de Realidade Aumentada que permitirá aos operadores consultar, junto da máquina identificada com uma anomalia, as instruções deixadas remotamente pelo técnico, auxiliando a resolução do problema sem necessidade de deslocação deste à fábrica.

\section{Tecnologias e Ferramentas}
A aplicação será desenvolvida nativamente para iOS, utilizando Swift como linguagem de programação e SwiftUI para a implementação da interface. Para as funcionalidades de Realidade Aumentada, será utilizada a framework ARKit, permitindo visualizar as instruções de manutenção junto das máquinas. A comunicação com o servidor central será realizada através de uma API, possibilitando o acesso a relatórios, vídeos, dados dos sensores e anotações do sistema AeroEye. O desenvolvimento será realizado numa máquina virtual com macOS, executada sobre o sistema operativo do computador de trabalho, seja Linux ou Windows.

% Inserir aqui o texto de contextualização do projeto.

\end{document}